\documentclass[CJK,13pt]{beamer}
\usepackage{CJKutf8}
\usepackage{beamerthemesplit}
\usetheme{Malmoe}
\useoutertheme[footline=authortitle]{miniframes}
\usepackage{amsmath}
\usepackage{amssymb}
\usepackage{graphicx}
\usepackage{eufrak}
\usepackage{color}
\usepackage{slashed}
\usepackage{simplewick}
\usepackage{tikz}
\usepackage{tcolorbox}
\usepackage{ulem}
\graphicspath{{../figures/}}
%%figures
\def\lfig#1#2{\includegraphics[width=#1 in]{#2}}
\def\tfig#1#2{\includegraphics[height=#1 in]{#2}}
\def\addfig#1#2{\begin{center}\includegraphics[width=#1 in]{#2}\end{center}}
\def\question{\includegraphics[width=0.3in]{why.jpg}\,}
\def\answer{\includegraphics[width=0.3in]{baozou_haha.png}\,}
\def\wulian{\includegraphics[width=0.18in]{emoji_wulian.jpg}}
\def\bigwulian{\includegraphics[width=0.35in]{emoji_wulian.jpg}}
\def\bye{\includegraphics[width=0.18in]{emoji_bye.jpg}}
\def\bigbye{\includegraphics[width=0.35in]{emoji_bye.jpg}}
\def\huaixiao{\includegraphics[width=0.18in]{emoji_huaixiao.jpg}}
\def\bighuaixiao{\includegraphics[width=0.35in]{emoji_huaixiao.jpg}}
\def\jianxiao{\includegraphics[width=0.18in]{emoji_jianxiao.jpg}}
\def\bigjianxiao{\includegraphics[width=0.35in]{emoji_jianxiao.jpg}}
\def\haoqi{\includegraphics[width=0.18in]{emoji_haoqi.jpg}}
%% colors
\def\blacktext#1{{\color{black}#1}}
\def\bluetext#1{{\color{blue}#1}}
\def\redtext#1{{\color{red}#1}}
\def\darkbluetext#1{{\color[rgb]{0,0.2,0.6}#1}}
\def\skybluetext#1{{\color[rgb]{0.2,0.7,1.}#1}}
\def\cyantext#1{{\color[rgb]{0.,0.5,0.5}#1}}
\def\greentext#1{{\color[rgb]{0,0.7,0.1}#1}}
\def\darkgray{\color[rgb]{0.2,0.2,0.2}}
\def\lightgray{\color[rgb]{0.6,0.6,0.6}}
\def\gray{\color[rgb]{0.4,0.4,0.4}}
\def\blue{\color{blue}}
\def\red{\color{red}}
\def\orange{\color[rgb]{1.,0.8,0.}}
\def\green{\color{green}}
\def\darkgreen{\color[rgb]{0,0.4,0.1}}
\def\darkblue{\color[rgb]{0,0.2,0.6}}
\def\skyblue{\color[rgb]{0.2,0.7,1.}}
%%control
\def\diag{\mathrm{diag}\,}
\def\heaviside{\,\mathrm{h}\,}
\def\bral{\left(\begin{array}{l}}
\def\brar{\end{array}\right)}
\def\brall{\left(\begin{array}{ll}}
\def\brarr{\end{array}\right)}
\def\bralll{\left(\begin{array}{lll}}
\def\brarrr{\end{array}\right)}
\def\branchl{\left\{\begin{array}{l}}
\def\branchr{\end{array}\right.}
\def\branchll{\left\{\begin{array}{ll}}
\def\branchrr{\end{array}\right.}
\def\branchlll{\left\{\begin{array}{lll}}
\def\branchrrr{\end{array}\right.}
\def\sfgamma#1{\,\Gamma\left( #1 \right)\,}
\def\be{\begin{equation}}
\def\ee{\nonumber\end{equation}}
\def\bea{\begin{eqnarray}}
\def\eea{\nonumber\end{eqnarray}}
\def\bch{\begin{CJK}{UTF8}{gbsn}}
\def\ech{\end{CJK}}
\def\bitem{\begin{itemize}}
\def\eitem{\end{itemize}}
\def\bcenter{\begin{center}}
\def\ecenter{\end{center}}
\def\bex{\begin{minipage}{0.2\textwidth}\includegraphics[width=0.6in]{jugelizi.png}\end{minipage}\begin{minipage}{0.76\textwidth}}
\def\eex{\end{minipage}}
\def\chtitle#1{\frametitle{\bch#1\ech}}
\def\bmat#1{\left(\begin{array}{#1}}
\def\emat{\end{array}\right)}
\def\bcase#1{\left\{\begin{array}{#1}}
\def\ecase{\end{array}\right.}
\def\bmini#1{\begin{minipage}{#1\textwidth}}
\def\emini{\end{minipage}}
\def\tbox#1{\begin{tcolorbox}#1\end{tcolorbox}}
\def\pfrac#1#2#3{\left(\frac{\partial #1}{\partial #2}\right)_{#3}}
\def\res#1#2{\,\mathrm{res}\,#1\left(#2\right)\,}
\def\newt#1#2{\left(\begin{array}{c}#1\\ #2\end{array}\right)}
\def\reof#1{\,\mathrm{Re}{\left(#1\right)\,}}
\def\imof#1{\,\mathrm{Im}{\left(#1\right)\,}}
\def\Arg#1{\,\mathrm{Arg}\,#1\,}
%%symbols
\def\intfull{\int_{-\infty}^\infty}
\def\inthalf{\int_0^\infty}
\def\bropt{\,(\ \ \ )}
\def\sech{\mathrm{sech}\,}
\def\csch{\mathrm{csch}\,}
\def\asinh{\mathrm{asinh}\,}
\def\acosh{\mathrm{acosh}\,}
\def\sone{$\star$}
\def\stwo{$\star\star$}
\def\sthree{$\star\star\star$}
\def\sfour{$\star\star\star\star$}
\def\sfive{$\star\star\star\star\star$}
\def\rint{{\int_\leftrightarrow}}
\def\roint{{\oint_\leftrightarrow}}
\def\stdHf{{\textit{\r H}_f}}
\def\deltaH{{\Delta \textit{\r H}}}
\def\ii{{\dot{\imath}}}
\def\skipline{{\vskip0.1in}}
\def\skiplines{{\vskip0.2in}}
\def\lagr{{\mathcal{L}}}
\def\hamil{{\mathcal{H}}}
\def\vecT{{\mathbf{T}}}
\def\vecN{{\mathbf{N}}}
\def\vecB{{\mathbf{B}}}
\def\vecv{{\mathbf{v}}}
\def\vecr{{\mathbf{r}}}
\def\vecf{{\mathbf{f}}}
\def\vecg{{\mathbf{g}}}
\def\vecupsilon{{\mathbf{\upsilon}}}
\def\vecu{{\mathbf{u}}}
\def\vecj{{\mathbf{j}}}
\def\vecx{{\mathbf{x}}}
\def\vecy{{\mathbf{y}}}
\def\vecz{{\mathbf{z}}}
\def\veck{{\mathbf{k}}}
\def\vecp{{\mathbf{p}}}
\def\vecn{{\mathbf{n}}}
\def\vecA{{\mathbf{A}}}
\def\vecP{{\mathbf{P}}}
\def\vecsigma{{\mathbf{\sigma}}}
\def\hatJn{{\hat{J_\vecn}}}
\def\hatJx{{\hat{J_x}}}
\def\hatJy{{\hat{J_y}}}
\def\hatJz{{\hat{J_z}}}
\def\hatj#1{\hat{J_{#1}}}
\def\hatphi{{\hat{\phi}}}
\def\hatq{{\hat{q}}}
\def\hatpi{{\hat{\pi}}}
\def\vel{\upsilon}
\def\Dint{{\mathcal{D}}}
\def\adag{{\hat{a}^\dagger}}
\def\bdag{{\hat{b}^\dagger}}
\def\cdag{{\hat{c}^\dagger}}
\def\ddag{{\hat{d}^\dagger}}
\def\hata{{\hat{a}}}
\def\hatb{{\hat{b}}}
\def\hatc{{\hat{c}}}
\def\hatd{{\hat{d}}}
\def\hatD{{\,\hat{D}}}
\def\hatN{{\hat{N}}}
\def\hatH{{\hat{H}}}
\def\hatp{{\hat{p}}}
\def\Fup{{F^{\mu\nu}}}
\def\Fdown{{F_{\mu\nu}}}
\def\newl{\nonumber \\}
\def\vece{\mathrm{e}}
\def\calM{{\mathcal{M}}}
\def\calT{{\mathcal{T}}}
\def\calR{{\mathcal{R}}}
\def\barpsi{\bar{\psi}}
\def\baru{\bar{u}}
\def\barv{\bar{\upsilon}}
\def\qeq{\stackrel{?}{=}}
\def\ftf{\stackrel{\mathcal{FT}}{\Longrightarrow}}
\def\ftb{\stackrel{\mathcal{FT}}{\Longleftarrow}}
\def\ftfb{\stackrel{\mathcal{FT}}{\Longleftrightarrow}}
\def\ltf{\stackrel{\mathcal{LT}}{\Longrightarrow}}
\def\ltb{\stackrel{\mathcal{LT}}{\Longleftarrow}}
\def\ltfb{\stackrel{\mathcal{LT}}{\Longleftrightarrow}}
\def\torder#1{\mathcal{T}\left(#1\right)}
\def\rorder#1{\mathcal{R}\left(#1\right)}
\def\contr#1#2{\contraction{}{#1}{}{#2}#1#2}
\def\trof#1{\mathrm{Tr}\left(#1\right)}
\def\trace{\mathrm{Tr}}
\def\comm#1{\ \ \ \left(\mathrm{used}\ #1\right)}
\def\tcomm#1{\ \ \ (\text{#1})}
\def\slp{\slashed{p}}
\def\slk{\slashed{k}}
\def\calp{{\mathfrak{p}}}
\def\veccalp{\mathbf{\mathfrak{p}}}
\def\Tthree{T_{\tiny \textcircled{3}}}
\def\pthree{p_{\tiny \textcircled{3}}}
\def\dbar{{\,\mathchar'26\mkern-12mu d}}
\def\erf{\mathrm{erf}}
\def\const{\mathrm{const.}}
\def\pheat{\pfrac p{\ln T}V}
\def\vheat{\pfrac V{\ln T}p}
%%units
\def\fdeg{{^\circ \mathrm{F}}}
\def\cdeg{^\circ \mathrm{C}}
\def\atm{\,\mathrm{atm}}
\def\angstrom{\,\text{\AA}}
\def\SIL{\,\mathrm{L}}
\def\SIkm{\,\mathrm{km}}
\def\SIyr{\,\mathrm{yr}}
\def\SIGyr{\,\mathrm{Gyr}}
\def\SIV{\,\mathrm{V}}
\def\SImV{\,\mathrm{mV}}
\def\SIeV{\,\mathrm{eV}}
\def\SIkeV{\,\mathrm{keV}}
\def\SIMeV{\,\mathrm{MeV}}
\def\SIGeV{\,\mathrm{GeV}}
\def\SIcal{\,\mathrm{cal}}
\def\SIkcal{\,\mathrm{kcal}}
\def\SImol{\,\mathrm{mol}}
\def\SIN{\,\mathrm{N}}
\def\SIHz{\,\mathrm{Hz}}
\def\SIm{\,\mathrm{m}}
\def\SIcm{\,\mathrm{cm}}
\def\SIfm{\,\mathrm{fm}}
\def\SImm{\,\mathrm{mm}}
\def\SInm{\,\mathrm{nm}}
\def\SImum{\,\mathrm{\mu m}}
\def\SIJ{\,\mathrm{J}}
\def\SIW{\,\mathrm{W}}
\def\SIkJ{\,\mathrm{kJ}}
\def\SIs{\,\mathrm{s}}
\def\SIkg{\,\mathrm{kg}}
\def\SIg{\,\mathrm{g}}
\def\SIK{\,\mathrm{K}}
\def\SImmHg{\,\mathrm{mmHg}}
\def\SIPa{\,\mathrm{Pa}}
%page
\def\secpage#1#2{\begin{frame}\bcenter{\bf \Huge #1} \ecenter \skipline \tbox{\bcenter #2 \ecenter}\end{frame}}
\def\append#1#2{\secpage{附录#1}{#2}}
\def\thinka#1{\begin{frame}\frametitle{思考题}\bcenter\lfig{0.4}{think0.jpg}  \skipline #1 \ecenter\end{frame}}
\def\thinkb#1{\begin{frame}\frametitle{思考题}\bcenter\lfig{0.5}{think1.jpg}  \skipline #1 \ecenter\end{frame}}
\def\thinkc#1{\begin{frame}\frametitle{思考题}\bcenter\lfig{0.5}{think2.jpg}  \skipline #1 \ecenter\end{frame}}
\def\thinkd#1{\begin{frame}\frametitle{思考题}\bcenter\lfig{0.5}{think3.jpg}  \skipline #1 \ecenter\end{frame}}
\def\thinke#1{\begin{frame}\frametitle{思考题}\bcenter\lfig{0.5}{think4.jpg}  \skipline #1 \ecenter\end{frame}}
\def\thinkf#1{\begin{frame}\frametitle{思考题}\bcenter\lfig{0.5}{think5.jpg}  \skipline #1 \ecenter\end{frame}}
\def\schw{ds^2 = \left(1-\frac{2GM}{r}\right)dt^2 - \left(1-\frac{2GM}{r}\right)^{-1} dr^2 - r^2\left(d\theta^2 + \sin^2\theta d\phi^2\right)}

\input{titlepage.tex}
  \date{}
  \begin{document}
  \bch
  \tpage{12}{Energy Momentum Tensor}


  \begin{frame}
    \frametitle{{\blue 物质} 的 {\blue 能量} {\blue 动量} {\blue 张量}}

    \addfig{1}{blackq.jpg}
    
    这四个词我都知道是什么意思，但合在一起……
  \end{frame}

  \secpage{关于符号的一点点疑惑}{学广义相对论最难搞清楚的不是指标，而是正负号}

  \begin{frame}
    上一讲我们给出了能量动量张量的定义：
    \tbox{$$T_{\mu\nu} = 2\frac{\delta \mathcal{L}_m}{\delta g^{\mu\nu}} -g_{\mu\nu}\mathcal{L}_m$$}
    或者等价地
    \tbox{$$T_{\mu\nu} = \frac{2}{\sqrt{-g}}\frac{\delta \left(\sqrt{-g}\mathcal{L}_m\right)}{\delta g^{\mu\nu}} $$}    
  \end{frame}


  \begin{frame}
    我们也可以以逆变的形式定义：
    \tbox{$$T^{\mu\nu} = -2\frac{\delta \mathcal{L}_m}{\delta g_{\mu\nu}} -g^{\mu\nu}\mathcal{L}_m$$}
    或者等价地
    \tbox{$$T^{\mu\nu} = -\frac{2}{\sqrt{-g}}\frac{\delta \left(\sqrt{-g}\mathcal{L}_m\right)}{\delta g_{\mu\nu}} $$}        
    这里的各种负号让你困惑了吗？
  \end{frame}
    
  \begin{frame}
    我们先来看一个恒等式：
    $$ g^{\mu\rho}g_{\rho\lambda} = g^{\mu}_{\ \lambda}$$
    两边取变分，右边是个常数，变分为零，所以
    $$ \delta g^{\mu\rho}g_{\rho\lambda}+g^{\mu\rho}\delta g_{\rho\lambda} = 0 $$
    两边同乘以 $g^{\nu\lambda}$，并把第二项移到方程右边，得到
    \tbox{$$ \delta g^{\mu\nu} = - g^{\mu\rho}g^{\nu\lambda}\delta g_{\rho\lambda}$$}
    类似地可以得到
    \tbox{$$ \delta g_{\mu\nu} = - g_{\mu\rho}g_{\nu\lambda}\delta g^{\rho\lambda}$$}    
  \end{frame}


  \begin{frame}
    也就是说：
    
    {\blue $\delta g_{\mu\nu}$ 和 $\delta g^{\mu\nu}$ 之间可以进行指标升降，但要多个负号。}

    于是有

    \bea
    \delta \mathcal{L}_m &=& f_{\mu\nu}\delta g^{\mu\nu} \newl
    &=& f^{\alpha\beta}g_{\alpha\mu}g_{\beta\nu}\delta g^{\mu\nu} \newl
    &=& -f^{\alpha\beta}\delta g_{\alpha\beta}
    \eea
    这里的 $f_{\mu\nu} = \frac{1}{2}\left(T_{\mu\nu}+g_{\mu\nu}\mathcal{L}_m\right)$。

    然后就有
    \bea
    T^{\alpha\beta} &=& 2 f^{\alpha\beta}  - g^{\alpha\beta}\mathcal{L}_m \newl
    &=& -2\frac{\delta\mathcal{L}_m}{\delta g_{\alpha\beta}} - g^{\alpha\beta}\mathcal{L}_m 
    \eea
    
  \end{frame}

  
  \begin{frame}
    \bcenter
    \lfig{1}{chulecai.jpg}
    
    这种数学推导除了让本喵更迷惑之外毫无用处……
    \ecenter
  \end{frame}


  \secpage{一个奇怪的标量}{长这么丑竟然也是……}

  \begin{frame}
    我们把四维坐标 $x^\mu$ 分离为 $t=x^0$ 和 $\vecx=(x^1,x^2,x^3)$ （注意这里的 $\vecx$ 只具有数学意义，不代表客观物理矢量），试说明
    设某静止质量为 $m$ 的粒子的世界线方程为 $\vecx = \vecy(t)$，四维动量为 $p^\mu=m\frac{dx^\mu}{ds}$ ($ds=\sqrt{g_{\mu\nu}dx^\mu dx^\nu}$是弧长参数)。试说明    
    $$\frac{\delta^{(3)}(\vecx-\vecy)}{p^0\sqrt{-g}}$$
    是标量，这里的 $\delta^{(3)}(\vecx-\vecy)$ 是（数学意义上的）三维Dirac函数。

    \addfig{1.5}{choubiaoliang.jpg}
  \end{frame}
  
  \secpage{一个粒子的$T_{\mu\nu}$}{$$ T^{\mu\nu} =  p^\mu p^\nu \frac{\delta^{(3)}\left(\vecx - \vecy\right)}{p^0\sqrt{-g}}$$}
  
  \begin{frame}
    \frametitle{粒子的 $T^{\mu\nu}$}    
    一个质量为 $m$ 的自由粒子，我们不妨称之为Bob，它有作用量
    $$ S_m = -m \int ds $$
    积分沿粒子的世界线进行，$s$ 为弧长参数。

    \skipline

    假设 Bob 的世界线用时间坐标 $t=x^0$ 描述为 $y^\mu(t)$，其四维动量为 $p^\mu = m\frac{dx^\mu}{dt}$，
    $$ S_m = -m \int \sqrt{g_{\mu\nu}\left(y(t)\right) \frac{dy^\mu}{dt}\frac{dy^\nu}{dt}}dt$$
    把粒子的四维动量记作 $p^\mu$，上式可以写成
    $$ S_m = - m\int \sqrt{g_{\mu\nu}\left(y(t)\right) p^\mu p^\nu}\frac{dt}{p^0}$$
  \end{frame}


  \begin{frame}
    \frametitle{粒子的 $T^{\mu\nu}$（续）}    
    利用三维 $\delta$ 函数，Bob的作用量可以写成：
    $$ S_m = - m\int \sqrt{g_{\mu\nu}\left(x\right)p^\mu p^\nu} \frac{\delta^{(3)}\left(\vecx - \vecy\right)}{p^0}d^3\vecx dt$$
    把这个式子河蟹一下：
    $$ S_m = - m\int \sqrt{g_{\mu\nu}\left(x\right)p^\mu p^\nu} \frac{\delta^{(3)}\left(\vecx - \vecy\right)}{p^0\sqrt{-g}} \sqrt{-g} d^4x$$
    于是 $\mathcal{L}_m = -  m\sqrt{g_{\mu\nu}\left(x\right)p^\mu p^\nu} \frac{\delta^{(3)}\left(\vecx - \vecy\right)}{p^0\sqrt{-g}}$.
    
    {\scriptsize 要说明一下，这里四维坐标系是给定的，作用量里采用的广义坐标是对应于每一个 $t$ 的 $y^\mu(t)$，以及对应于每一个四维点 $x$ 的 $g_{\mu\nu}(x)$。我们曾经固定$g_{\mu\nu}$，变动 $y^\mu(t)$，推导了粒子的运动方程；现在我们固定 $y^\mu(t)$，变动 $g_{\mu\nu}$，来推导爱因斯坦方程里的 $T^{\mu\nu}$（如果你要推导完整的爱因斯坦方程，当然需要把依赖于 $g_{\mu\nu}$ 的时空作用量加进来一起做变分）。}    
  \end{frame}

  \begin{frame}
    \frametitle{粒子的 $T^{\mu\nu}$（续）}
    最后，

    \bea
    T^{\mu\nu} &=& -\frac{2}{\sqrt{-g}}\frac{\delta\left(\sqrt{-g}\mathcal{L}_m\right)}{\delta g_{\mu\nu}} \newl
    &=& m\frac{p^\mu p^\nu}{ \sqrt{g_{\mu\nu}\left(x\right)p^\mu p^\nu}} \frac{\delta^{(3)}\left(\vecx - \vecy\right)}{p^0\sqrt{-g}}
    \eea
    由于已经对 $g_{\mu\nu}$ 取完了变分，我们可以取 $g_{\mu\nu}$ 为“该取的值”了，也就是有 $g_{\mu\nu}p^\mu p^\nu = p^\mu p_\mu = m^2$ ，{\blue 最后我们得到单粒子的能量动量张量：
   $$ T^{\mu\nu} =  p^\mu p^\nu \frac{\delta^{(3)}\left(\vecx - \vecy\right)}{p^0\sqrt{-g}}$$}
  \end{frame}


  \begin{frame}
    \frametitle{粒子的 $T^{\mu\nu}$（续）}
    如果某四维点 $x$ 附近，有一堆粒子作{\blue 统计上各向同性的随机运动}，则
    $$T^{0i} \propto \sum p^i = 0$$
    且
    $$T^{ij} \propto \sum p^i p^j \propto \delta^{ij}$$
    即能量动量张量是对角矩阵，且 $T^{11}=T^{22}=T^{33}$。


    \skipline
    
    和平直时空的非相对论理想气体比较，我们可以把 $T^{00}$ 叫做能量密度，把 $T^{11}=T^{22}=T^{33}$ 叫做压强。

  \end{frame}
  

  \secpage{理想流体的 $T^{\mu\nu}$}{$$T^{\mu\nu} = (p+\rho)u^\mu u^\nu - p g^{\mu\nu}$$}

  \begin{frame}
    \frametitle{理想流体的 $T^{\mu\nu}$}
    “理想流体”的每一个四维时空点 $x$ 都有一个“宏观速度” $u^\mu(x)$，和流体宏观速度共动的观测者看到附近的流体是各向同性的。流体的实质当然是一堆有微观随机运动的粒子，所以流体的共动参照系里流体的能量动量张量为：
    \be
    \left.\tilde{T}^{\mu\nu}\right\vert_{\rm comoving} = \begin{pmatrix}
      \rho & 0 & 0 & 0 \\
      0 & p & 0 & 0 \\
      0 & 0 & p & 0 \\
      0 & 0 & 0 & p      
    \end{pmatrix}
    \ee
    这里的 $\rho$ 是能量密度，$p$ 是压强（均指共动参考系里观测到的）。
  \end{frame}

  
  
  \begin{frame}
    \frametitle{理想流体的 $T^{\mu\nu}$}    
    转换到原先的参照系，我们先给结论： 理想流体的能量动量张量是
     \tbox{$$T^{\mu\nu} =  (\rho+p)u^\mu u^\nu - p g^{\mu\nu}$$}
    {\scriptsize 如果你采用 $-+++$ 度规符号习惯，上面的右边第二项的减号要变成加号。}

    \skiplines

     如何证明这个 $T^{\mu\nu}$ 在共动参考系里具有 $\diag(\rho, p, p, p)$ 的形式呢？下一讲我们将详细讨论观测量的问题。
  \end{frame}


  
  
  
\ech
\end{document}




  
